\section{Chapter 1}

\begin{quotex}
In 1959 \textbf{Julius Evola} published a translation and commentary on the \emph{Tao Te Ching} under the title \emph{Il Libro del Principio e della sua Azione} The Book of the Principle and its Action . His interpretation is based on his knowledge of Tradition.

Many years prior he had written an alternative translation that was based more on his studies of German Idealism than on Tradition.
\end{quotex}


\begin{quotex}
\begin{verse}
The Tao that can be named\\
Is not the eternal Tao\\
The Name that can be pronounced\\
Is not the eternal Name\\
(As the) Nameless it is the principle of Heaven and Earth\\
With the Name that is: determined as Heaven and Earth it is the origin of the myriad particular beings\\
So: the one who is detached\\
Perceives the Mysterious Essence\\
Which is obscured by desire\\
The gaze is stopped by the limit he sees only the phenomena of appearances of the Principle \\
Now of the two the Nameable and the Unnamable, being and non-being \\
One essence, only the name is different\\
Their identity is mysterious\\
It is the unfathomable depth\\
Beyond the threshold of the ultimate mystery
\end{verse}
\end{quotex}


The two aspects of the Principle are indicated here, the one transcendent (unnameable, without form --- equivalent to non-Being, the Emptiness, the unmoved), the other immanent (nameable, with form --- equivalent to Being, fullness, the movable). Other Taoist designations: The ``Former Heaven'' (\emph{hsien t'ien}) and the ``Later Heaven'' (\emph{hou t'ien}). The first determination of the Tao (called also ``the great Mutation'') is Heaven-and-Earth which are cosmic symbols of the Yang and the Yin. This Dyad produces all the modifications, therefore the myriad particular beings, of their paths and of their destinies (``The Great Flux'', the ``current of forms'').

The unmanifested and the manifested, the formless and the formal, the Former Heaven and the Later Heaven, are not distinct temporally (as if at a certain moment something similar to a ``creation'' intervened), but in logico-metaphysical terms. Transcendence is immanent, the fullness coexists with the emptiness, non-being is the source of being: identity, that constitutes the ultimate mystery of Taoist realization.

See \emph{Lieh Tzu}:

\begin{quotex}
The chain of productions and transformations is interrupted, the producer and the transformer producing and transforming without it --- the producer is unmovable, the transformer comes and goes. And the motionless and the mobile endure always. ... Analyzing the production of the universe, the opening up of the sensible from the non-sensible, the seed of the calm generative action of Heaven and Earth, the ancient Sages distinguished these stages: the Great Mutation, the Great Origin, the Great Commencement, the Great Flux.
\end{quotex}


\paragraph{Chinese text and literal translation}


Chapter 1 (\begin{CJK*}{UTF8}{bsmi}第一章\end{CJK*})

\begin{table}[h]
\centering
\begin{tabular}{p{.35\textwidth}p{.5\textwidth}}

\begin{minipage}[t]{.35\textwidth}
\begin{CJK*}{UTF8}{bsmi}
道可道，非恆道；

名可名，非恆名。

無名天地之始；

有名萬物之母。

故，

恆無，欲也，以觀其妙；

恆有，欲也，以觀其徼。

此兩者同出而異名，

同謂之玄。

玄之又玄，眾妙之門。
\end{CJK*}
\end{minipage}
 &
\begin{minipage}[t]{.5\textwidth}
The Dao that can be stated, is not the eternal Dao;

The name that can be named, is not the eternal name.

The nameless is the origin of heaven and earth;

The named is the mother of the myriad things.

So,

By constantly having no desire one views its wonders;

By constantly having desire , one views its limits.

These two have the same origin, but they differ in name;

Both are called Mystery.

One Mystery plus another Mystery, is the source of all wonders.
\end{minipage}

\end{tabular}
\end{table}


\flrightit{Posted on 2020-03-12 by Lao Tzu}

